%%%%%%%%%%%%%%%%%%%%%%%%%%%%%%%%%%%%%%%%
%      Fichier maître (Main.tex)      %
%%%%%%%%%%%%%%%%%%%%%%%%%%%%%%%%%%%%%%%%
% Dorian Depriester, 2014

\usepackage{lmodern}
\usepackage{ae,aecompl}
\usepackage{eso-pic}	% Nécessaire pour mettre des images en arrière plan
\usepackage{array}		% Permet d'écrite 'THESE' de haut en bas
\usepackage[letterspace=-60]{microtype}
%\usepackage{microtype}
%police, si pas de base dans la feuille de style 
\renewcommand{\rmdefault}{pnc} 
\renewcommand{\sfdefault}{pag} 
\renewcommand{\ttdefault}{pcr}
%%%%%%%%%%%%%%%%%%%%%%%%%%%%%%%%%%%%%%%%
%    Page de garde (Pagedegarde.tex)   %
%%%%%%%%%%%%%%%%%%%%%%%%%%%%%%%%%%%%%%%%
% Dorian Depriester, 2014

\makeatletter
\def\@ecole{école}
\newcommand{\ecole}[1]{
  \def\@ecole{#1}
}

\def\@specialite{Spécialité}
\newcommand{\specialite}[1]{
  \def\@specialite{#1}
}

\def\@ED{\'{E}cole Doctorale}
\newcommand{\ED}[1]{
  \def\@ED{#1}
}

\def\@doctorat{Doctorat}
\newcommand{\doctorat}[1]{
  \def\@doctorat{#1}
}

\def\@adresse{Adresse}
\newcommand{\adresse}[1]{
  \def\@adresse{#1}
}

\def\@directeur{directeur}
\newcommand{\directeur}[1]{
  \def\@directeur{#1}
}

\def\@encadrant{encadrant}
\newcommand{\encadrant}[1]{
  \def\@encadrant{#1}
}
\def\@jurya{}{}{}
\newcommand{\jurya}[3]{
  \def\@jurya{#1,	& #2	& #3\\}
}
\def\@juryb{}{}{}
\newcommand{\juryb}[3]{
  \def\@juryb{#1,	& #2	& #3\\}
}
\def\@juryc{}{}{}
\newcommand{\juryc}[3]{
  \def\@juryc{#1,	& #2	& #3\\}
}
\def\@juryd{}{}{}
\newcommand{\juryd}[3]{
  \def\@juryd{#1,	& #2	& #3\\}
}
\def\@jurye{}{}{}
\newcommand{\jurye}[3]{
  \def\@jurye{#1,	& #2	& #3\\}
}
\def\@juryf{}{}{}
\newcommand{\juryf}[3]{
  \def\@juryf{#1,	& #2	& #3\\}
}
\def\@juryg{}{}{}
\newcommand{\juryg}[3]{
  \def\@juryg{#1,	& #2	& #3\\}
}
\def\@juryh{}{}{}
\newcommand{\juryh}[3]{
  \def\@juryh{#1,	& #2	& #3\\}
}
\def\@juryi{}{}{}
\newcommand{\juryi}[3]{
  \def\@juryi{#1,	& #2	& #3\\}
}
\makeatother

\newcommand\BackgroundPic{%
	\put(0,0){%
		\parbox[b][\paperheight]{\paperwidth}{%
			\hspace{5em}
			\includegraphics[height = 0.45\paperheight]{Illustrations/bordure}%
			\vfill
		}
	}
}
\newcommand\EtiquetteThese{%
%Ici c'est le positionnement de l'encadré en bas à droite
	\put(-20,65){% c'est le point bas droite
		\parbox[t][\paperheight]{\paperwidth}{%
			\hfill
			\colorbox{black}{		
				\begin{minipage}[b]{3em}
				\vspace{0.3cm}        %ça c'est bon
					\centering\Huge\textcolor{white}{T\\H\\\`{E}\\S\\E\\}
					\vspace{0.3cm} %ça c'est bon
				\end{minipage}
			}
		}
	}
}

\makeatletter
\newcommand{\pagedegarde}{
 \newgeometry{top=0.8cm, bottom=1cm, left=2cm, right=1cm}
 \AddToShipoutPicture*{\BackgroundPic}
 \AddToShipoutPicture*{\EtiquetteThese}
   \begin{titlepage}
   \centering
   \hspace{-0.5cm}
       \includegraphics[width=0.35\textwidth]{Illustrations/paristech}
       \hfill
       \includegraphics[width=0.35\textwidth]{Illustrations/ensam}\\
     \vspace{-0.2cm}
     \hspace{7.4cm}
 \small 2018-ENAM-0019
      \vspace{0.2cm}
           \hspace{-7.4cm}

      {\Large \@ED}\\
    \vspace{1cm}
      {\huge 
      	{\bfseries \@doctorat}\\
              %\Large \color{red}(Mémoire provisoire)\color{black}\\
                  \vspace{.5cm} %a mettre si non provisoire
    \vspace{.5cm}
      \huge	TH\`{E}SE}\\
    \vspace{0.6cm}
   		{\bfseries pour obtenir le grade de docteur délivré par}\\
    \vspace{1cm}
    	{\huge\bfseries \@ecole}\\
    \vspace{0.5cm}
    	{\Large{\bfseries Spécialité ``\@specialite''}}\\
    \vspace{1cm}
    	\textit{présentée et soutenue publiquement par}\\
    \vspace{0.5cm}
    	{\Large {\bfseries \@author}} \\
    \vspace{0.5cm}
    	le \@date \\
    \vfill
       {\huge \color[rgb]{0,0,0} \bfseries{{\setlength{\baselineskip}{1\baselineskip} \@title \par}}}
    \vfill
        Directeur de thèse : {\bfseries \@directeur}\\
        Co-encadrants de thèse : {\bfseries \@encadrant}\\
    \vfill
    \hspace{-1cm}
	\begin{tabular}{>{\bfseries}lll}
		\large Jury\\
            \@juryf
		\@juryd
		\@jurye
		\@jurya
		\@juryb
		\@juryc
		\@juryg
		\@juryh
		\@juryi
	\end{tabular}
	\vfill
	
	\@adresse
   \end{titlepage}




 \restoregeometry  
}

\author{Antoine \textsc{Lassagne}}
\title{Retour sensoriel du toucher dans des environnements virtuels}
\ED{{\fontfamily{bch}\selectfont \'{E}cole doctorale \no 432 : Sciences des Métiers de l'Ingénieur}}
\doctorat{{\fontfamily{bch}\selectfont Doctorat ParisTech}}
\specialite{{\fontfamily{bch}\selectfont Informatique -- traitement du signal}}
\directeur{Andras \textsc{Kemeny}}
\encadrant{Frederic \textsc{Merienne} et Javier \textsc{Posselt}}
\date{28 mai 2018}
\jurya{M. Andras Kemeny}{Professeur, Arts \& Métiers}{Examinateur}
\juryb{M. Frédéric Merienne}{Professeur, Arts \& Métiers}{Examinateur}
\juryc{M. Javier Posselt}{Ingénieur, Renaut S.A.S.}{Examinateur}
\juryd{M. Jean-Christophe Popieul}{Professeur, Univ. de Valenciennes}{Rapporteur}
\jurye{M. Franck Mars}{Chargé de Recherche, LS2N}{Rapporteur}
\juryf{M. Christophe Chaillou}{Professeur, Université de Lille}{Président}
\ecole{\lsstyle {\fontfamily{ppl}\selectfont ~~l'\'{E}cole Nationale Supérieure d'Arts et Métiers}}
\adresse{
\textbf{Arts et Métiers ParisTech}\\Laboratoire d'Immersion Virtuelle -- Institut Image
}

% Et tout le reste du document
%\end{document}